\documentclass[11pt,letterpaper]{article}

% ============================================================
% PACKAGES
% ============================================================

\usepackage[utf8]{inputenc}
\usepackage[T1]{fontenc}
\usepackage{lmodern}

\usepackage[
    letterpaper,
    margin=0.8in,
    top=0.8in,
    bottom=0.8in
]{geometry}

\usepackage{amsmath,amssymb}
\usepackage{siunitx}

\sisetup{
    per-mode=symbol,
    range-phrase={--},
    range-units=single
}

\usepackage{
    graphicx,
    booktabs,
    tabularx,
    array,
    ragged2e,
    float,
    caption,
    enumitem
}

\usepackage{microtype}

\usepackage{xcolor}
\usepackage[most]{tcolorbox}

\usepackage{fancyhdr}
\usepackage{titlesec}
\usepackage{lastpage}
\usepackage{needspace}

\usepackage[
    colorlinks=true,
    linkcolor=HIFENavy,
    citecolor=HIFETeal,
    urlcolor=HIFETeal,
    breaklinks=true,
    pdfstartview=FitH
]{hyperref}

\urlstyle{same}

% ============================================================
% COLORS
% ============================================================

\definecolor{HIFENavy}{RGB}{20,45,70}
\definecolor{HIFETeal}{RGB}{0,105,115}
\definecolor{HIFELight}{RGB}{240,246,248}
\definecolor{HIFEGray}{RGB}{100,105,110}

% ============================================================
% DOCUMENT MACROS
% ============================================================

\newcommand{\TODO}[1]{\textcolor{red}{[#1]}}

\newcommand{\ExhibitOneFile}{IMG_1484.jpeg}
\newcommand{\ExhibitTwoFile}{IMG_1483.jpeg}

% ============================================================
% PAGE STYLE
% ============================================================

\setlength{\headheight}{14pt}

\pagestyle{fancy}
\fancyhf{}

\lhead{%
    \small\textcolor{HIFENavy}{%
        HIFE: Hierarchical Integrated Flow Electrode
    }%
}

\rhead{%
    \small\textcolor{HIFEGray}{%
        Storage Design STEP Prize --- Phase 1
    }%
}

\cfoot{%
    \small\textcolor{HIFEGray}{%
        \thepage\ of \pageref{LastPage}
    }%
}

\renewcommand{\headrulewidth}{0.4pt}
\renewcommand{\footrulewidth}{0pt}

% ============================================================
% SECTION FORMATTING
% ============================================================

\titleformat
    {\section}
    {\Large\bfseries\color{HIFENavy}}
    {}
    {0em}
    {}

\titleformat
    {\subsection}
    {\large\bfseries\color{HIFETeal}}
    {}
    {0em}
    {}

\titleformat
    {\subsubsection}
    {\normalsize\bfseries\color{HIFENavy}}
    {}
    {0em}
    {}

\titlespacing*{\section}{0pt}{1.2em}{0.5em}
\titlespacing*{\subsection}{0pt}{1em}{0.4em}
\titlespacing*{\subsubsection}{0pt}{0.8em}{0.3em}

% ============================================================
% GENERAL TYPOGRAPHY
% ============================================================

\setlength{\parindent}{0pt}
\setlength{\parskip}{0.5em}

\setlength{\tabcolsep}{4pt}
\renewcommand{\arraystretch}{1.15}

\emergencystretch=2em

\captionsetup{
    font=small,
    labelfont=bf,
    justification=raggedright,
    singlelinecheck=false
}

\setlist[itemize]{
    leftmargin=1.6em,
    itemsep=2pt,
    topsep=3pt
}

% ============================================================
% EXHIBIT HEADING
% ============================================================

\newcommand{\exhibitheading}[1]{%
    \needspace{8\baselineskip}%
    \begin{tcolorbox}[
        colback=HIFENavy,
        colframe=HIFENavy,
        boxrule=0pt,
        arc=1.5mm,
        left=4mm,
        right=4mm,
        top=1.5mm,
        bottom=1.5mm
    ]
    \color{white}\bfseries #1
    \end{tcolorbox}%
}

% ============================================================
% TABLE COLUMN TYPE
% ============================================================

\newcolumntype{L}[1]{%
    >{\RaggedRight\arraybackslash}p{#1}%
}

% ============================================================
% WORD COUNTS
% ============================================================

\newcommand{\BGcount}{\TODO{recount; maximum 500}}
\newcommand{\TNcount}{2{,}819}

% ============================================================
% DOCUMENT
% ============================================================

\begin{document}

% ============================================================
% COVER PAGE
% ============================================================

\thispagestyle{empty}

\vspace*{0.15in}

{\color{HIFENavy}\rule{\textwidth}{1.2pt}}

\vspace{0.15in}

\begin{center}

    {\LARGE\bfseries
    HIFE: Hierarchical Integrated Flow Electrode\par}

    \vspace{0.08in}

    {\large\bfseries
    A Manufacturable All-Iron Aqueous Redox-Flow Battery Electrode\par}

    \vspace{0.08in}

    {\large
    Storage Design Strategies to Ease Production (STEP) Prize --- Phase 1\par}

\end{center}

\vspace{0.15in}

{\color{HIFETeal}\rule{\textwidth}{0.8pt}}

\vspace{0.18in}

\noindent
{\large\bfseries\color{HIFENavy}
Project Information\par}

\vspace{0.45em}

\noindent
\textbf{Project title:}\par
HIFE: Hierarchical Integrated Flow Electrode

\vspace{0.30em}

\noindent
\textbf{Competitor name:}\par
Umar Tabbsum

\vspace{0.30em}

\noindent
\textbf{Location:}\par
Dublin, Ohio

\vspace{0.30em}

\noindent
\textbf{Short project description:}\par
HIFE is an electrode architecture for aqueous all-iron redox-flow batteries
that assigns electrochemical surface area, permeability, and flow
distribution to separate, individually controllable length scales. The
architecture is designed to specify hydraulic performance and manufacturing
tolerance together, using abundant iron chemistry and manufacturing
processes intended to be compatible with existing sheet and roll-to-roll
equipment. This Phase~1 submission identifies key production and
supply-chain challenges, proposes design-for-manufacturability responses,
and defines validation gates for subsequent development.

\vspace{0.30em}

\noindent
\textbf{Introduction video (90 s):}\par
\url{https://youtu.be/5L1X5Yqj6oA}

\vspace{0.30em}

\noindent
\textbf{Main point of contact:}\par
Umar Tabbsum\\
\href{mailto:umartabbs@yahoo.com}{umartabbs@yahoo.com}

\vspace{0.30em}

\noindent
\textbf{Evidence status:}\par
Analytical and computational screening only; no experimental or CFD
validation has been performed to date.

\vspace{0.30em}

\noindent
\textbf{Submission date:}\par
October 03, 2026

\vfill

{\color{HIFENavy}\rule{\textwidth}{1.2pt}}

\newpage

% ============================================================
% COMPETITOR BACKGROUND
% ============================================================

\section{Competitor Background}

I hold dual B.S. degrees in Computer Engineering Technology and Electronics
Engineering Technology from DeVry University (2005), followed by a
ten-year career in test engineering and accredited laboratory operations
(2006--2016) and subsequent experience in regulated production and
operations.

\textbf{Test engineering and quality systems.}
From 2010 to 2016 at Intertek (Grand Rapids, MI), I advanced from Test
Technician~III to EMC Engineer, Engineering Team Lead, and ultimately Senior
Project EMC Engineer. I co-built and helped manage an
ISO/IEC~17025/A2LA-accredited EMC laboratory, working with the EMC Laboratory
Manager on laboratory development, accreditation activities, and technical
operations. I developed LabVIEW-based test automation and custom IEEE~488.2
(GPIB) instrument drivers and managed multi-year client programs serving
automotive, aerospace, and defense customers, including responsibility for
project budgets and P\&L performance. I received three Intertek Excellence
Awards, including recognition associated with an ISO/IEC~17025/A2LA
accreditation audit.

\textbf{Electrified transport and batteries.}
My work included EMC testing of electric-vehicle charging equipment
[\textit{insert exact product/equipment and role}] and EMC testing of a
Nissan Leaf traction-battery system. I coordinated with the Automotive EMC
Test Engineer to troubleshoot test anomalies using automotive component-level
EMC methods, including bulk current injection (BCI) in accordance with
\textbf{ISO 11452-4}, \textit{Road vehicles --- Component test methods for
electrical disturbances from narrowband radiated electromagnetic energy ---
Part 4: Bulk current injection (BCI)}. This work provided practical exposure
to the electromagnetic compatibility of high-voltage automotive systems,
battery-pack architecture, charging behavior, and troubleshooting of
abnormal test responses. First-generation Nissan Leaf traction batteries
used laminated lithium-ion cells with a manganese-spinel-based cathode,
providing additional practical context for cell construction, charging
behavior, and battery-system failure modes.

\textbf{Materials chemistry and emissions.}
At Intertek I reviewed volatile-organic-compound (VOC) chamber testing with
the professor who led that laboratory, to learn how methods such as
ANSI/BIFMA~M7.1 and CDPH Standard Method~V1.2 identify volatile-emission
problems in furniture and bedding materials. This experience taught me how
to trace a material's chemistry to a measurable compliance risk, a
discipline HIFE applies to electrolytes, binders, and pore formers.

\textbf{Aerospace programs.}
At GE Aviation Digital Systems (2006--2010), I operated large-scale complex
test systems supporting more than eight concurrent defense and commercial
programs.

\textbf{Professional scope.}
I have not worked professionally in electrochemistry or porous-electrode
fabrication. The screening model, figures, and plan in this submission are
my own work in MATLAB. Fabrication, electrochemical characterization, and CFD
would be conducted with partners (see Commercialization), and the validation
gates ensure that partner data, rather than assumptions, determine each
development step.

% ============================================================
% SOLUTION DESCRIPTION
% ============================================================

\section{Solution Description}

\subsection{What HIFE Is}

HIFE is a hierarchical electrode architecture for aqueous all-iron
redox-flow batteries intended for grid-scale stationary long-duration
energy storage. The design separates functions that are often coupled in a
conventional porous electrode by assigning electrochemical surface area,
permeability, and flow distribution to distinct but connected length
scales.

The architecture contains four coupled scales:

\begin{itemize}
    \item \textbf{Microscale fibrous region:} provides high accessible
    surface area for electrochemical reactions.

    \item \textbf{Mesoscale porous network:} establishes controlled
    permeability and tortuosity.

    \item \textbf{Engineered hydraulic pathways:} provide preferential
    transport routes to reduce maldistribution and excessive pressure
    drop.

    \item \textbf{Cell-level distribution plate:} distributes electrolyte
    across the active electrode area.
\end{itemize}

HIFE therefore treats hydraulic resistance and manufacturability as design
variables rather than secondary consequences of electrode fabrication.

\subsection{Novelty and Technical Differentiation}

The novelty claim is directed to the electrode and flow-path architecture,
not to the underlying all-iron electrolyte chemistry.

Iron-based flow batteries are being developed for utility-scale applications,
including multi-megawatt and multi-megawatt-hour systems. HIFE addresses a
different layer of the system: integration of graded porous structure and
engineered hydraulic pathways into a manufacturable electrode component
co-designed with the cell flow architecture.

The proposed architecture is intended to provide a controllable relationship
between:

\begin{enumerate}
    \item electrochemically active area;
    \item hydraulic permeability;
    \item flow distribution;
    \item compression tolerance;
    \item manufacturing variation; and
    \item resulting pressure drop and pumping requirement.
\end{enumerate}

No HIFE electrode has yet been fabricated. Consequently, pore uniformity,
channel dimensional tolerance, compression response, fabrication yield,
material sourcing, and electrochemical performance remain open engineering
questions.

\subsection{Evidence and Current Development Status}

The present evidence base is analytical and computational screening. The
model is intentionally first-order and is not presented as experimental
validation or computational-fluid-dynamics validation.

For a porous electrode, the screening pressure-drop relationship is

\begin{equation}
\Delta P
=
\frac{\mu L v}{k},
\label{eq:pressure-drop}
\end{equation}

where $\mu$ is dynamic viscosity, $L$ is electrode thickness, $v$ is
superficial velocity, and $k$ is permeability.

Permeability is estimated using a Kozeny--Carman-type relationship:

\begin{equation}
k =
\frac{d_f^2\varepsilon^3}
{180(1-\varepsilon)^2}.
\label{eq:permeability}
\end{equation}

The nominal screening assumptions are:

\begin{itemize}
    \item dynamic viscosity:
    $\mu=\SI{1.5}{mPa.s}$;
    \item porosity:
    $\varepsilon=0.90$;
    \item fiber diameter:
    $d_f=\SI{10}{\micro\meter}$;
    \item electrode thickness:
    $L=\SI{3}{mm}$;
    \item superficial velocity:
    $v=\SI{0.01}{m.s^{-1}}$.
\end{itemize}

A 20,000-realization Monte Carlo analysis was performed using random seed
7. The uncertainty ranges were viscosity $\pm10\%$, porosity
$0.88$--$0.92$, fiber diameter $\pm10\%$, and thickness $\pm10\%$.

The resulting pressure-drop distribution had a median of
\SI{1.107}{kPa}, with a 5th--95th percentile range of
\SIrange{0.658}{1.774}{kPa}.

Pearson sensitivity correlations with pressure drop were approximately:

\begin{itemize}
    \item porosity: $r=-0.88$;
    \item fiber diameter: $r=-0.38$;
    \item viscosity: $r=+0.19$;
    \item thickness: $r=+0.18$.
\end{itemize}

These results identify porosity as a high-leverage design variable in the
first-order model.

The model excludes compression, entrance effects, gas evolution, clogging,
and other phenomena that could materially affect an operating cell.

% ============================================================
% HYDRAULIC DESIGN
% ============================================================

\section{Hydraulic and Manufacturing Design}

\subsection{Pressure Drop and Pumping Requirement}

The associated pumping-power relationship is

\begin{equation}
P_{\mathrm{pump}}
=
\frac{\Delta P\,Q}
{\eta_{\mathrm{pump}}},
\label{eq:pump-power}
\end{equation}

where $Q$ is volumetric flow rate and $\eta_{\mathrm{pump}}$ is pump
efficiency.

The purpose of the present calculation is not to establish a final pump
rating. Instead, it establishes a screening framework linking electrode
geometry to hydraulic burden.

\subsection{Porosity Sensitivity}

The pressure-drop relationship implies

\begin{equation}
\Delta P
\propto
\frac{(1-\varepsilon)^2}{\varepsilon^3}.
\label{eq:porosity-dependence}
\end{equation}

The logarithmic sensitivity is

\begin{equation}
\frac{\mathrm{d}\ln\Delta P}
{\mathrm{d}\varepsilon}
=
-\frac{2}{1-\varepsilon}
-\frac{3}{\varepsilon}.
\label{eq:porosity-sensitivity}
\end{equation}

At $\varepsilon=0.90$, this derivative is approximately $-23$. This
illustrates why relatively small changes in porosity can produce substantial
changes in predicted pressure drop.

An illustrative $\pm20\%$ pressure-drop target corresponds locally to a
porosity change on the order of $\pm0.01$ around $\varepsilon=0.90$ under
the screening relationship. The actual working tolerance is a validation
question and would be established through Gate~2 rather than assumed.

% ============================================================
% COST AND MANUFACTURING
% ============================================================

\section{Cost and Manufacturing Strategy}

\subsection{Accepted-Cost Framework}

A useful manufacturing metric is accepted cost rather than nominal material
cost:

\begin{equation}
C_{\mathrm{acc}}
=
\frac{
C_{\mathrm{mat}}
+
C_{\mathrm{lab}}
+
C_{\mathrm{ovh}}
}{Y},
\label{eq:accepted-cost}
\end{equation}

where $C_{\mathrm{mat}}$ is material cost, $C_{\mathrm{lab}}$ is direct
labor cost, $C_{\mathrm{ovh}}$ is allocated overhead, and $Y$ is accepted
yield.

This formulation explicitly captures the effect of manufacturing yield on
the effective cost of usable electrodes.

\subsection{Design-for-Manufacturability}

The development strategy is to avoid requiring an entirely new manufacturing
ecosystem. Candidate processes should be evaluated for compatibility with
existing sheet-processing and roll-to-roll equipment where technically
appropriate.

The principal manufacturing variables are:

\begin{itemize}
    \item fiber diameter;
    \item fiber loading;
    \item porosity;
    \item thickness;
    \item binder content;
    \item channel dimensions;
    \item compression;
    \item dimensional uniformity; and
    \item production yield.
\end{itemize}

Manufacturing specifications should be established from measured process
capability rather than imposed as arbitrary targets.

% ============================================================
% SUPPLY CHAIN
% ============================================================

\section{Supply-Chain Strategy}

The supply-chain strategy emphasizes multiple qualified sources for critical
inputs and minimizes dependence on specialized single-source materials.

HIFE avoids relying on vanadium, cobalt, nickel, and lithium as active redox
materials. The all-iron chemistry provides an opportunity to build the
electrode system around more widely available materials while retaining
attention to electrode-specific supply constraints.

Supplier concentration can be monitored using the Herfindahl--Hirschman
Index:

\begin{equation}
\mathrm{HHI}
=
\sum_i s_i^2,
\label{eq:hhi}
\end{equation}

where $s_i$ is supplier $i$'s share of the relevant supply.

The objective is not to eliminate all supplier concentration but to identify
materials where qualification of a second source materially reduces
production risk.

% ============================================================
% COMMERCIALIZATION
% ============================================================

\section{Commercialization and Development Plan}

The proposed development sequence is:

\begin{itemize}
    \item \textbf{0--6 months:} fabricate screening coupons and establish
    baseline material and hydraulic properties.

    \item \textbf{6--12 months:} validate pressure-drop behavior,
    compression response, dimensional uniformity, and initial
    electrochemical performance.

    \item \textbf{12--24 months:} fabricate larger-area electrodes,
    demonstrate repeatability, and validate manufacturing yield.

    \item \textbf{24--36 months:} integrate the electrode into a
    representative flow-battery stack and evaluate long-duration
    operation.
\end{itemize}

The principal commercial pathway is therefore staged validation rather than
immediate scale-up.

% ============================================================
% AI DISCLOSURE
% ============================================================

\section{AI-Assisted Development Disclosure}

\TODO{Confirm the exact AI tools and uses before submission.}

Generative AI tools were used to assist with document organization, language
editing, LaTeX formatting, and review of the presentation against the
competition requirements. The competitor developed and/or approved the
technical content, screening model, figures, engineering assumptions, and
development plan and remains responsible for the accuracy of all statements.

% ============================================================
% EXHIBITS
% ============================================================

\section{Exhibits}

% ------------------------------------------------------------
% EXHIBIT 1
% ------------------------------------------------------------

\needspace{12\baselineskip}

\exhibitheading{Exhibit 1 --- Conceptual HIFE Architecture}

\begin{samepage}
\begin{figure}[H]
\centering
\includegraphics[
    width=\textwidth,
    keepaspectratio
]{\ExhibitOneFile}

\caption{\textbf{Conceptual HIFE architecture.}
Layers shown: (1)~cell-level distribution plate,
(2)~hydraulic pathways, (3)~mesoscale porous network,
(4)~fibrous electrode, and (5)~membrane.
Conceptual schematic; not CFD, experimental imaging, or a fabricated
prototype.}

\label{fig:hife-architecture}
\end{figure}
\end{samepage}

% ------------------------------------------------------------
% EXHIBIT 2
% ------------------------------------------------------------

\needspace{12\baselineskip}

\exhibitheading{Exhibit 2 --- Hydraulic Screening}

\begin{samepage}
\begin{figure}[H]
\centering
\includegraphics[
    width=\textwidth,
    keepaspectratio
]{\ExhibitTwoFile}

\caption{\textbf{Hydraulic screening.}
(A)~Nominal inputs and scope.
(B)~Porous-electrode pressure drop, 20,000 Monte Carlo realizations at
\SI{0.01}{m.s^{-1}}: median \SI{1107}{Pa}, 5th--95th percentile
\SIrange{658.4}{1774}{Pa}.
(C)~Baseline versus channel-assisted surrogate
(4,000 realizations), each normalized to its own median; relative response
only, not an absolute pressure-drop or pump-power result.
(D)~Pearson correlations with $\log_{10}\Delta P$
(3,000 realizations).
Analytical screening; seed 7; no experimental or CFD data.}

\label{fig:hydraulic-screening}
\end{figure}
\end{samepage}

% ------------------------------------------------------------
% EXHIBIT 3
% ------------------------------------------------------------

\needspace{16\baselineskip}

\exhibitheading{Exhibit 3 --- Manufacturing Challenges, Design Responses, and Metrics}

\begin{samepage}
\small
\begin{center}
\begin{tabularx}{\textwidth}{
    L{0.27\textwidth}
    L{0.36\textwidth}
    L{0.29\textwidth}
}
\toprule
\textbf{Manufacturing challenge}
&
\textbf{Design response}
&
\textbf{Validation metric}
\\
\midrule

Porosity variation
&
Control fiber loading, compaction, and process conditions
&
Porosity distribution; pressure-drop repeatability
\\

Fiber-diameter variation
&
Specify qualified fiber-size distribution and incoming inspection
&
Diameter distribution; permeability variation
\\

Thickness variation
&
Define process-control limits and in-line dimensional inspection
&
Thickness Cpk; pressure-drop variation
\\

Channel dimensional tolerance
&
Use manufacturable channel geometry with controlled land/channel ratio
&
Channel width/depth tolerance; flow distribution
\\

Compression sensitivity
&
Characterize electrode response under representative compression
&
Pressure drop versus compression; thickness recovery
\\

Yield loss
&
Track defect modes and establish process capability before scale-up
&
Accepted yield; defect Pareto; accepted cost
\\

Material sourcing
&
Qualify multiple suppliers for critical inputs
&
Supplier qualification; lead time; HHI
\\

Scale-up repeatability
&
Transfer validated process windows to larger-area manufacturing
&
Lot-to-lot variation; hydraulic and electrochemical repeatability
\\

\bottomrule
\end{tabularx}
\end{center}
\normalsize
\end{samepage}

% ------------------------------------------------------------
% EXHIBIT 4
% ------------------------------------------------------------

\needspace{16\baselineskip}

\exhibitheading{Exhibit 4 --- Supply-Chain Dependency Map}

\begin{samepage}
\small
\begin{center}
\begin{tabularx}{\textwidth}{
    L{0.24\textwidth}
    L{0.28\textwidth}
    L{0.23\textwidth}
    L{0.17\textwidth}
}
\toprule
\textbf{Input}
&
\textbf{Primary dependency}
&
\textbf{Potential mitigation}
&
\textbf{Monitoring metric}
\\
\midrule

Conductive fiber / carbon material
&
Qualified material supplier and consistent morphology
&
Second-source qualification
&
Lot-to-lot property variation
\\

Binder
&
Chemical compatibility and supply continuity
&
Dual sourcing and incoming QC
&
Viscosity; solids content; lead time
\\

Porogen / pore-forming input
&
Particle-size distribution and availability
&
Alternative grades and suppliers
&
Particle-size distribution; supplier concentration
\\

Current collector
&
Electrical conductivity and dimensional consistency
&
Qualified alternate source
&
Resistance; dimensional tolerance
\\

Membrane
&
Chemical compatibility and area availability
&
Multiple qualified suppliers
&
Area yield; lead time
\\

Fabrication equipment
&
Coating, drying, calendaring, and patterning capability
&
Use existing industrial equipment where feasible
&
Throughput; uptime; process capability
\\

\bottomrule
\end{tabularx}
\end{center}
\normalsize
\end{samepage}

% ------------------------------------------------------------
% EXHIBIT 5
% ------------------------------------------------------------

\needspace{16\baselineskip}

\exhibitheading{Exhibit 5 --- Validation Gates and Principal Risks}

\begin{samepage}
\small
\begin{center}
\begin{tabularx}{\textwidth}{
    L{0.10\textwidth}
    L{0.39\textwidth}
    L{0.41\textwidth}
}
\toprule
\textbf{Gate}
&
\textbf{Validation objective}
&
\textbf{Principal risk addressed}
\\
\midrule

G1
&
Fabricate repeatable HIFE coupons with controlled thickness and porosity
&
Architecture may not be manufacturable at the required dimensional
tolerance
\\

G2
&
Measure pressure drop and establish the acceptable porosity/process window
&
Analytical permeability model may not represent fabricated electrodes
\\

G3
&
Characterize compression response and hydraulic stability
&
Stack compression may alter permeability and flow distribution
\\

G4
&
Demonstrate electrochemical performance and cycling stability
&
Hydraulic optimization may not preserve electrochemical activity
\\

G5
&
Demonstrate repeatability and production yield at larger area
&
Small-coupon performance may not translate to manufacturing scale
\\

G6
&
Integrate validated electrodes into a representative flow-battery stack
&
System-level interactions may dominate cell performance
\\

\bottomrule
\end{tabularx}
\end{center}
\normalsize
\end{samepage}

% ============================================================
% REFERENCES
% ============================================================

\section{References}

\begin{enumerate}[leftmargin=2em,itemsep=4pt]

\item U.S. Department of Energy,
\textit{Achieving the Promise of Low-Cost Long Duration Energy Storage},
August 2024.
\href{
https://www.energy.gov/sites/default/files/2024-08/Achieving\%20the\%20Promise\%20of\%20Low-Cost\%20Long\%20Duration\%20Energy\%20Storage_FINAL_08052024.pdf
}{
DOE Long Duration Energy Storage Report
}.

\item ISO 11452-4,
\textit{Road vehicles --- Component test methods for electrical disturbances
from narrowband radiated electromagnetic energy --- Part 4: Bulk current
injection (BCI)}.

\item ANSI/BIFMA M7.1,
\textit{Standard Test Method for Determining VOC Emissions from Office
Furniture Systems, Components and Seating}.

\item California Department of Public Health,
\textit{Standard Method for the Testing and Evaluation of Volatile Organic
Chemical Emissions from Indoor Sources Using Environmental Chambers},
Standard Method V1.2.

\item Kozeny, J.,
\textit{Über kapillare Leitung des Wassers im Boden},
Sitzungsberichte der Akademie der Wissenschaften in Wien, 1927.

\item Carman, P. C.,
\textit{Fluid Flow through Granular Beds},
Transactions of the Institution of Chemical Engineers, 1937.

\item Forward Edge-AI / HeroX,
\textit{Blaise xTech Competition},
competition materials and Phase~I requirements.

\item U.S. Department of Energy,
\textit{Long Duration Energy Storage for the Electricity Grid},
program and technology materials.

\item ESS Tech, Inc.,
publicly available materials describing iron-flow battery technology and
utility-scale deployment.

\item Publicly available literature on aqueous iron redox-flow batteries,
porous electrodes, hydraulic resistance, and flow distribution.

\end{enumerate}

% ============================================================
% END DOCUMENT
% ============================================================

\end{document}
